% =====================================================================
%  BÖLÜM 6: Dönüşümler ve Dağılım Teorisi
%  Kaynak: Hogg & Craig; Casella & Berger; Feller Vol.II; Klenke
% =====================================================================
\chapter{Dönüşümler ve Dağılım Teorisi}
\label{chp:donusumler}

\bolumalinti{%
  Bir rasgele değişkenin dönüşümü, dağılım uzayındaki bir yolculuktur.%
}{George Casella \& Roger Berger, \textit{Statistical Inference}}

% ---- Kazanımlar (TYMM) ----
\begin{kazanimlar}
Bu bölümün sonunda öğrenci:
\begin{itemize}[leftmargin=1.8em]
  \item Tekil ve çok boyutlu dönüşümlerde Jacobian formülünü ustalıkla uygulayabilir;
        birebir olmayan dönüşümleri parçalara bölerek ele alabilir.
  \item MÜF ve karakteristik fonksiyon yöntemiyle dağılım tanımlayabilir;
        bu yöntemin Jacobian yöntemine göre avantajlarını açıklayabilir.
  \item Evrişim (konvolüsyon) integralini hesaplayabilir; Gama, Normal ve Poisson
        toplamlarının dağılımını türetebilir.
  \item Sıra istatistiklerini tanımlayabilir; tekil ve ortak PDF'lerini türetip
        hesaplayabilir; örneklem minimum, maksimum ve medyanını bulabilir.
  \item Delta yöntemini tek ve çok boyutlu durumda uygulayabilir.
\end{itemize}
\end{kazanimlar}

% ---- Motivasyon ----
\begin{motivasyon}
Bölüm~\ref{chp:rasgele_degiskenler}'de dönüşüm yönteminin temellerini gördük.
Bu bölümde araç kutusunu genişletiyoruz. Çoğu zaman Jacobian hesabı yerine
MÜF ya da KF yöntemi çok daha hızlı sonuç verir: $X+Y$'nin dağılımını bulmak için
evrişim integralini hesaplamak yerine $M_{X+Y}(t) = M_X(t)M_Y(t)$ özdeşliğini
kullanmak yeterlidir. Sıra istatistikleri ise güvenilirlik mühendisliği ve
çevresel istatistiğin temel aracıdır: "en zayıf halka ne zaman kopar?"

\medskip
\textbf{Köprü.} Bölüm~\ref{chp:rasgele_degiskenler}'deki Jacobian ve evrişim,
Bölüm~\ref{chp:beklenti}'deki MÜF/KF bu bölümde birleştiriliyor.
Bölüm~\ref{chp:merkezi_limit}'teki Delta yöntemi bu bölümün doğal devamıdır.
\end{motivasyon}

% =====================================================================
\section{Tekil Dönüşümler: Derinlemesine}
\label{sec:tekil_donusum}
% =====================================================================

Bölüm~\ref{sec:rd_donusum}'da birebir dönüşüm formülünü ve genel monoton-olmayan
durumu gördük. Burada bu tekniği pekiştiriyor, sınır durumlarını ve pratikte
sık karşılaşılan kalıpları inceliyoruz.

\begin{teorem}[Tekil Dönüşüm — Kapsamlı Form]
\label{thm:tekil_donusum}
$X$ sürekli rasgele değişken, PDF $f_X$. $g : \mathcal{X} \to \mathcal{Y}$
($\mathcal{X}$ ve $\mathcal{Y}$ açık aralıklar birleşimi). $\mathcal{X} = \bigsqcup_k A_k$
ayrışımında $g|_{A_k}$ her parça üzerinde birebir, türevlenebilir olsun.
$h_k = (g|_{A_k})^{-1}$ ters fonksiyon. O zaman $Y = g(X)$'in PDF'si:
\[
  f_Y(y) = \sum_k f_X(h_k(y)) \cdot |h_k'(y)| \cdot \mathbf{1}_{y \in g(A_k)}.
\]
\end{teorem}

\begin{ornek}[$Y = |X|$, $X \sim \Normal(0,1)$]
$g(x) = |x|$: $A_1 = (-\infty, 0)$, $A_2 = (0, \infty)$.
$h_1(y) = -y$, $h_2(y) = y$; $|h_1'| = |h_2'| = 1$.
$f_Y(y) = \phi(-y) + \phi(y) = 2\phi(y)$, $y > 0$ (yarı-normal dağılım).
$P(|Z| > 1.96) = 2(1 - \Phi(1.96)) \approx 0.05$ — iki-kuyruklu normal tablosu.
\end{ornek}

\begin{ornek}[$Y = X^2$, $X$ simetrik]
$X$ simetrik (çift PDF) olsun. $A_1 = (-\infty,0)$, $A_2 = (0,\infty)$.
$h_1(y) = -\sqrt{y}$, $h_2(y) = \sqrt{y}$; $|h_{1,2}'(y)| = 1/(2\sqrt{y})$.
$f_Y(y) = [f_X(\sqrt{y}) + f_X(-\sqrt{y})]/(2\sqrt{y}) = f_X(\sqrt{y})/\sqrt{y}$ (simetri).
$Z \sim \Normal(0,1)$: $f_{Z^2}(y) = \phi(\sqrt{y})/\sqrt{y} = (2\pi y)^{-1/2}e^{-y/2}$
$= \Gama(1/2, 1/2)$ PDF'si $= \chi^2_1$ PDF'si.
\end{ornek}

\begin{ispatiskele}[Quantile Dönüşümü — İspat İskeleti]
$U \sim \Uniform(0,1)$, $F$ herhangi KDF, $X = F^{-1}(U)$.
\iskeleadim{1}{$P(X \leq x) = P(F^{-1}(U) \leq x)$ ifadesini yazın.}
\iskeleadim{2}{$F^{-1}(u) \leq x \Leftrightarrow u \leq F(x)$ eşdeğerliğini kullanın ($F$ sağdan sürekli).}
\iskeleadim{3}{$P(U \leq F(x)) = \bosluk{F(x)}$ (uniform KDF).}
\iskeleadim{4}{Sonuç: $F_X(x) = F(x)$, yani $X \sim F$.}
\end{ispatiskele}

% =====================================================================
\section{Çok Boyutlu Dönüşümler}
\label{sec:cok_boyut_donusum}
% =====================================================================

Bölüm~\ref{subsec:donusum_jacobian}'da Jacobian formülünü verdik.
Burada sistematik uygulama tekniklerini ve önemli örnekleri işliyoruz.

\begin{onerme}[Jacobian Hesabında Adım Adım Yöntem]
\label{prop:jacobian_adim}
$(X_1, \ldots, X_n) \to (Y_1, \ldots, Y_n)$ dönüşümü için:
\begin{enumerate}[(1)]
  \item Ters dönüşümü bulun: $x_i = h_i(y_1,\ldots,y_n)$.
  \item Jacobian matrisini oluşturun: $J_{ij} = \partial h_i/\partial y_j$.
  \item $|\det J|$'yi hesaplayın.
  \item Tanım bölgesini dönüştürün: $\mathcal{X} \to \mathcal{Y}$.
  \item $f_Y(\mathbf{y}) = f_X(h(\mathbf{y})) \cdot |\det J|$ ile PDF'yi yazın.
\end{enumerate}
\end{onerme}

\begin{ornek}[Beta-Gama İlişkisi — Tam Türetme]
\label{ex:beta_gama_turetme}
$U \sim \Gama(\alpha, 1)$, $V \sim \Gama(\beta, 1)$ bağımsız.
$X = U/(U+V)$, $W = U+V$ dönüşümü.

\textbf{Ters:} $u = xw$, $v = w(1-x)$; $0 < x < 1$, $w > 0$.

\textbf{Jacobian:}
$J = \begin{pmatrix} w & x \\ -w & 1-x \end{pmatrix}$; $|\det J| = w(1-x) + xw = w$.

\textbf{Ortak PDF:}
\begin{align*}
  f_{X,W}(x,w) &= f_U(xw)\,f_V(w(1-x))\cdot w \\
  &= \frac{(xw)^{\alpha-1}e^{-xw}}{\Gamma(\alpha)} \cdot \frac{(w(1-x))^{\beta-1}e^{-w(1-x)}}{\Gamma(\beta)} \cdot w \\
  &= \frac{x^{\alpha-1}(1-x)^{\beta-1}}{B(\alpha,\beta)} \cdot \frac{w^{\alpha+\beta-1}e^{-w}}{\Gamma(\alpha+\beta)}.
\end{align*}

$X$ ve $W$ bağımsız; $X \sim \Beta(\alpha,\beta)$, $W \sim \Gama(\alpha+\beta, 1)$.
\end{ornek}

\begin{ornek}[$t$ Dağılımının Türetilmesi]
\label{ex:t_turetme}
$Z \sim \Normal(0,1)$, $V \sim \chi^2_n$ bağımsız. $T = Z/\sqrt{V/n}$, $W = V$.

\textbf{Ters:} $z = t\sqrt{w/n}$, $v = w$.
\textbf{Jacobian:} $|\det J| = \sqrt{w/n}$.

\textbf{Ortak PDF:}
$f_{T,W}(t,w) = \phi(t\sqrt{w/n}) \cdot f_{\chi^2_n}(w) \cdot \sqrt{w/n}$.

$W$'ye göre integre et ($w > 0$):
\[
  f_T(t) = \frac{\Gamma((n+1)/2)}{\sqrt{n\pi}\,\Gamma(n/2)} \left(1 + \frac{t^2}{n}\right)^{-(n+1)/2}.
\]
(Gama integrali $\int_0^\infty w^{(n-1)/2} e^{-w(1+t^2/n)/2} dw = \Gamma((n+1)/2) \cdot [2/(1+t^2/n)]^{(n+1)/2}$
kullanılır.)
\end{ornek}

\begin{ornek}[$F$ Dağılımının Türetilmesi]
$U \sim \chi^2_m$, $V \sim \chi^2_n$ bağımsız. $W = (U/m)/(V/n)$ için:
$f_W(w) = \frac{\Gamma((m+n)/2)}{\Gamma(m/2)\Gamma(n/2)} \left(\frac{m}{n}\right)^{m/2}
w^{m/2-1}\left(1 + \frac{m}{n}w\right)^{-(m+n)/2}$, $w > 0$.
Bu $F(m,n)$ dağılımının PDF'sidir.
\end{ornek}

% =====================================================================
\section{MÜF ve KF Yöntemiyle Dağılım Belirleme}
\label{sec:muf_yontemi}
% =====================================================================

\begin{motivasyon}[Jacobian Yerine MÜF]
$X+Y$'nin PDF'sini bulmak için ya evrişim integrali ya da MÜF çarpımı kullanılır.
MÜF yöntemi genellikle çok daha hızlıdır: toplam PDF'si yerine MÜF bulunur,
bu standart bir dağılıma karşılık geliyorsa PDF'yi yazmaya gerek kalmaz.
\end{motivasyon}

\begin{teorem}[MÜF ile Dağılım Tanımlama]
\label{thm:muf_tanim}
$Y = g(X_1, \ldots, X_n)$ ve $M_Y(t) = E[e^{tY}]$ hesaplanabiliyorsa:
$M_Y(t)$ bilinen bir dağılımın MÜF'üne eşit ise $Y$ o dağılımı izler
(MÜF tek türlü dağılım belirler, Teorem~\ref{thm:muf_tek}).
\end{teorem}

\begin{ornek}[Normal Toplamı — MÜF ile]
$X_i \sim \Normal(\mu_i, \sigma_i^2)$ bağımsız ($i=1,\ldots,n$), $a_i \in \mathbb{R}$.
\[
  M_{\sum a_i X_i}(t) = \prod_i M_{X_i}(a_i t) = \prod_i e^{\mu_i a_i t + \sigma_i^2 a_i^2 t^2/2}
  = \exp\!\left(\sum a_i\mu_i \cdot t + \frac{t^2}{2}\sum a_i^2\sigma_i^2\right).
\]
Bu $\Normal(\sum a_i\mu_i, \sum a_i^2\sigma_i^2)$ MÜF'üdür. Jacobian veya evrişim gerekmez.
\end{ornek}

\begin{ornek}[Gama Toplamı — MÜF ile]
$X_i \sim \Gama(\alpha_i, \beta)$ bağımsız ($\beta$ aynı):
$M_{\sum X_i}(t) = \prod (\beta/(\beta-t))^{\alpha_i} = (\beta/(\beta-t))^{\sum\alpha_i}$.
$\Rightarrow$ $\sum X_i \sim \Gama(\sum\alpha_i, \beta)$.

Özel durum: $Z_i \sim \Normal(0,1)$ bağımsız $\Rightarrow$ $Z_i^2 \sim \chi^2_1 = \Gama(1/2,1/2)$
$\Rightarrow$ $\sum_{i=1}^n Z_i^2 \sim \Gama(n/2,1/2) = \chi^2_n$.
\end{ornek}

\begin{ornek}[Poisson Toplamı — KF ile]
$X_i \sim \Pois(\lambda_i)$ bağımsız:
$\varphi_{\sum X_i}(t) = \prod e^{\lambda_i(e^{it}-1)} = \exp\!\bigl((\sum\lambda_i)(e^{it}-1)\bigr)$.
$\Rightarrow$ $\sum X_i \sim \Pois(\sum\lambda_i)$.
\end{ornek}

\begin{teorem}[KF ile Dağılımla Yakınsama]
\label{thm:kf_yaklinsama}
$\varphi_{X_n}(t) \to \varphi_X(t)$ her $t$'de ise $X_n \xrightarrow{d} X$
(Lévy süreklilik teoremi; ispat Bölüm~\ref{chp:merkezi_limit}'te kullanılacak).
\end{teorem}

% =====================================================================
\section{Evrişim}
\label{sec:evrisim}
% =====================================================================

Bölüm~\ref{subsec:evrisim}'de evrişim formülünü tanımladık. Burada hesaplama
tekniklerini ve önemli örnekleri derinleştiriyoruz.

\begin{teorem}[Evrişim Özellikleri]
\label{thm:evrisim_ozellik}
$f, g, h$ PDF'ler olmak üzere:
\begin{enumerate}[(i)]
  \item \textbf{Değişmelilik:} $f * g = g * f$.
  \item \textbf{Birleşmelilik:} $(f * g) * h = f * (g * h)$.
  \item \textbf{KF çarpımı:} $\varphi_{X+Y}(t) = \varphi_X(t)\,\varphi_Y(t)$.
  \item \textbf{MÜF çarpımı:} $M_{X+Y}(t) = M_X(t)\,M_Y(t)$ (bağımsız $X,Y$).
\end{enumerate}
\end{teorem}

\begin{ornek}[Üçgen Dağılım — Evrişim]
$X, Y \sim \Uniform(0,1)$ bağımsız, $Z = X + Y$:
\[
  f_Z(z) = \int_0^1 \mathbf{1}_{[0,1]}(z-y)\,dy
  = \begin{cases} z & 0 \leq z \leq 1 \\ 2-z & 1 < z \leq 2 \end{cases}.
\]
Üçgen dağılım (Bölüm~\ref{sec:buyuk_sayilar}'daki Al.7'nin açılımı).
$n$ bağımsız $\Uniform(0,1)$ toplamı: Irwin-Hall dağılımı.
$n \to \infty$: Merkezi Limit Teoremi'ne hazırlık.
\end{ornek}

\begin{ornek}[Normal Evrişim — İntegral Yöntemi]
$X \sim \Normal(\mu_1, \sigma_1^2)$, $Y \sim \Normal(\mu_2, \sigma_2^2)$ bağımsız.
Evrişim:
\[
  f_{X+Y}(z) = \int_{-\infty}^{+\infty} \frac{1}{\sigma_1\sqrt{2\pi}}e^{-(z-y-\mu_1)^2/(2\sigma_1^2)}
  \cdot \frac{1}{\sigma_2\sqrt{2\pi}}e^{-(y-\mu_2)^2/(2\sigma_2^2)} dy.
\]
Üstel ifade birleştirilip $y$'deki kareyi tamamlama:
$(z-y-\mu_1)^2/\sigma_1^2 + (y-\mu_2)^2/\sigma_2^2$'nin $y$ üzerinde integral
$\sigma_1^2+\sigma_2^2$ varyanslı Gauss integraline döner.
Sonuç: $X+Y \sim \Normal(\mu_1+\mu_2, \sigma_1^2+\sigma_2^2)$.
(MÜF yöntemi bu hesabı beş satıra indirger.)
\end{ornek}

\begin{onerme}[Cauchy Toplamı]
$X, Y \sim \mathrm{Cauchy}(0,1)$ bağımsız $\Rightarrow$ $X+Y \sim \mathrm{Cauchy}(0,2)$.
\end{onerme}

\begin{proof}
KF: $\varphi_{X+Y}(t) = e^{-|t|} \cdot e^{-|t|} = e^{-2|t|}$ — bu $\mathrm{Cauchy}(0,2)$'nin KF'sidir.
\end{proof}

\begin{kritiksorgu}
Cauchy toplamı neden $\mathrm{Cauchy}(0,2)$, "beklentiden" elde ettiğimiz $\mathrm{Cauchy}(0,1)$ değil?
Cauchy'nin beklentisi mevcut değildir — standartlaştırma ($n$'e bölme) dağılımı stabilize etmez.
Büyük Sayılar Kanunu çöker; Merkezi Limit Teoremi uygulanmaz.
\end{kritiksorgu}

% =====================================================================
\section{Sıra İstatistikleri}
\label{sec:sira_istatistikleri}
% =====================================================================

\begin{tanim}[Sıra İstatistikleri]
\label{def:sira_istatistik}
$X_1, \ldots, X_n$ i.i.d., ortak PDF $f$, KDF $F$.
Küçükten büyüğe sıralanmış değerler:
\[
  X_{(1)} \leq X_{(2)} \leq \cdots \leq X_{(n)}.
\]
$X_{(k)}$: $k$'ıncı \textbf{sıra istatistiği} (order statistic).
$X_{(1)} = \min$, $X_{(n)} = \max$, $X_{(\lceil n/2 \rceil)}$: medyan.
\end{tanim}

\begin{teorem}[Tekil Sıra İstatistiğinin PDF'si]
\label{thm:sira_pdf}
$k$'ıncı sıra istatistiğinin PDF'si:
\[
  f_{X_{(k)}}(x) = \frac{n!}{(k-1)!(n-k)!} [F(x)]^{k-1}[1-F(x)]^{n-k} f(x).
\]
\end{teorem}

\begin{proof}
Kombinatoryal argüman: $X_{(k)} \in (x, x+dx)$ olması için
$k-1$ gözlem $(-\infty, x)$'de, 1 gözlem $(x, x+dx)$'de, $n-k$ gözlem $(x+dx, \infty)$'da olmalı.
Çok terimli katsayı $\frac{n!}{(k-1)!\,1!\,(n-k)!}$ ile olasılık:
$[F(x)]^{k-1} \cdot f(x)\,dx \cdot [1-F(x)]^{n-k}$.
$dx$'e bölerek PDF elde edilir.
\end{proof}

\begin{ornek}[Minimum ve Maksimum]
\textbf{Minimum} ($k=1$):
$f_{X_{(1)}}(x) = n[1-F(x)]^{n-1}f(x)$.
$F_{X_{(1)}}(x) = 1 - [1-F(x)]^n$.

\textbf{Maksimum} ($k=n$):
$f_{X_{(n)}}(x) = n[F(x)]^{n-1}f(x)$.
$F_{X_{(n)}}(x) = [F(x)]^n$.

Bu formüller hemen doğrulanabilir:
$P(X_{(1)} > x) = P(\text{tüm } X_i > x) = [1-F(x)]^n$ (bağımsızlık).
$P(X_{(n)} \leq x) = P(\text{tüm } X_i \leq x) = [F(x)]^n$.
\end{ornek}

\begin{teorem}[Ortak PDF]
\label{thm:sira_ortak_pdf}
$(X_{(1)}, \ldots, X_{(n)})$'nin ortak PDF'si:
\[
  f_{X_{(1)},\ldots,X_{(n)}}(x_1, \ldots, x_n) = n! \prod_{i=1}^n f(x_i), \quad x_1 < x_2 < \cdots < x_n.
\]
\end{teorem}

\begin{proof}
$n!$ adet permütasyon; her biri $(x_1, \ldots, x_n)$ sıralamasını $x_{\sigma(1)} < \cdots < x_{\sigma(n)}$
ile üretir. Bağımsızlıktan her permütasyonun katkısı $\prod f(x_i)$; toplam $n! \prod f(x_i)$.
\end{proof}

\begin{teorem}[İkili Marjinal PDF]
\label{thm:sira_ikili}
$j < k$ için $(X_{(j)}, X_{(k)})$'nın ortak PDF'si:
\[
  f_{X_{(j)},X_{(k)}}(u,v) = \frac{n!}{(j-1)!(k-j-1)!(n-k)!}
  [F(u)]^{j-1}[F(v)-F(u)]^{k-j-1}[1-F(v)]^{n-k}f(u)f(v),
\]
$u < v$.
\end{teorem}

\begin{ornek}[$\Uniform(0,1)$ Sıra İstatistikleri]
$X_1, \ldots, X_n \sim \Uniform(0,1)$ i.i.d.; $f(x) = 1$, $F(x) = x$ ($0<x<1$).

$k$'ıncı sıra istatistiği:
$f_{X_{(k)}}(x) = \frac{n!}{(k-1)!(n-k)!} x^{k-1}(1-x)^{n-k} = \frac{x^{k-1}(1-x)^{n-k}}{B(k, n-k+1)}$.

$X_{(k)} \sim \Beta(k, n-k+1)$.
$E[X_{(k)}] = k/(n+1)$: $k$'ıncı sıra istatistiğinin beklentisi $[0,1]$'i eşit böler.
\end{ornek}

\begin{ornek}[Üstel Sıra İstatistikleri — Yarışma Süresi]
$X_i \sim \Exp(\lambda)$ bağımsız (cihaz ömürleri). $X_{(1)} = \min X_i$:
$P(X_{(1)} > t) = [P(X_i > t)]^n = e^{-n\lambda t}$.
$X_{(1)} \sim \Exp(n\lambda)$: minimum de üstel, oran $n$ kat büyük.

Sıra arası boşluklar: $D_k = X_{(k)} - X_{(k-1)}$ bağımsız,
$D_k \sim \Exp((n-k+1)\lambda)$ (eksponensiyal gecikme özelliği).
Bu, güvenilirlik mühendisliğinde yaygın kullanılan "rekabetçi hayatta kalma" modelidir.
\end{ornek}

\begin{hataanalizi}
"$X_{(k)}$ ile $X_{(j)}$ bağımsızdır" — yanlış.
Ortak PDF $f_{X_{(j)},X_{(k)}}(u,v) \neq f_{X_{(j)}}(u) f_{X_{(k)}}(v)$
(bölge $u < v$ koşulu bağımlılık yaratır).
Sıra istatistikleri $\mathbf{bağımlıdır}$; ikili ve yüksek dereceli dağılımlar ayrıca hesaplanmalıdır.
\end{hataanalizi}

\begin{mikroozet}
$f_{X_{(k)}}(x) = \frac{n!}{(k-1)!(n-k)!}F^{k-1}(1-F)^{n-k}f$.
Uniform: $X_{(k)} \sim \Beta(k, n-k+1)$.
Üstel: minimum de üstel; aralıklar bağımsız üstel.
Sıra istatistikleri genel olarak bağımlıdır.
\end{mikroozet}

% =====================================================================
\section{Delta Yöntemi}
\label{sec:delta_yontemi}
% =====================================================================

\begin{motivasyon}[Neden Delta Yöntemi?]
$\bar{X}_n \xrightarrow{d} \Normal(\mu, \sigma^2/n)$ — bu Merkezi Limit Teoremi'nin (bir sonraki
bölüm) sonucudur. Peki $g(\bar{X}_n)$ nasıl dağılır? $g$ doğrusal değilse doğrudan
dönüşüm uygulanamaz; ancak Taylor açılımı bu boşluğu kapatır.
\end{motivasyon}

\begin{teorem}[Delta Yöntemi — Tek Boyutlu]
\label{thm:delta_yontemi}
$\sqrt{n}(X_n - \mu) \xrightarrow{d} \Normal(0, \sigma^2)$ ve $g$ $\mu$'da türevlenebilir,
$g'(\mu) \neq 0$ ise:
\[
  \sqrt{n}(g(X_n) - g(\mu)) \xrightarrow{d} \Normal(0,\, \sigma^2 [g'(\mu)]^2).
\]
\end{teorem}

\begin{proof}
Taylor: $g(X_n) = g(\mu) + g'(\mu)(X_n-\mu) + R_n$ ($R_n = o(X_n-\mu)$).
$\sqrt{n}(g(X_n)-g(\mu)) = g'(\mu)\sqrt{n}(X_n-\mu) + \sqrt{n}R_n$.
$X_n \xrightarrow{P} \mu$ ($\sqrt{n}$-normalite bunu içerir) $\Rightarrow$ $\sqrt{n}R_n \xrightarrow{P} 0$.
Slutsky (Teorem~\ref{thm:slutsky}): $g'(\mu)\sqrt{n}(X_n-\mu) \xrightarrow{d} g'(\mu)\Normal(0,\sigma^2)
= \Normal(0, \sigma^2[g'(\mu)]^2)$.
\end{proof}

\begin{ornek}[Delta Yöntemi — Logaritma]
$X_i \sim \mathrm{Exp}(\lambda)$, $n$ büyük. MLT: $\sqrt{n}(\bar{X}_n - 1/\lambda) \xrightarrow{d} \Normal(0, 1/\lambda^2)$.
$g(x) = \ln x$, $g'(x) = 1/x$, $g'(1/\lambda) = \lambda$.
Delta yöntemi:
\[
  \sqrt{n}(\ln\bar{X}_n - \ln(1/\lambda)) \xrightarrow{d} \Normal(0, \lambda^2 \cdot 1/\lambda^2) = \Normal(0,1).
\]
\end{ornek}

\begin{teorem}[Çok Boyutlu Delta Yöntemi]
\label{thm:delta_cok}
$\sqrt{n}(\mathbf{X}_n - \boldsymbol{\mu}) \xrightarrow{d} \Normal_k(\mathbf{0}, \boldsymbol{\Sigma})$
ve $g : \mathbb{R}^k \to \mathbb{R}$ $\boldsymbol{\mu}$'da türevlenebilir ise:
\[
  \sqrt{n}(g(\mathbf{X}_n) - g(\boldsymbol{\mu})) \xrightarrow{d}
  \Normal\!\left(0,\, \nabla g(\boldsymbol{\mu})^\top \boldsymbol{\Sigma}\, \nabla g(\boldsymbol{\mu})\right).
\]
\end{teorem}

\begin{ornek}[Oran Tahmincisi]
$(X_i, Y_i)$ i.i.d., $E[X] = \mu_X$, $E[Y] = \mu_Y > 0$. $R_n = \bar{X}_n/\bar{Y}_n$.
$g(x,y) = x/y$; $\nabla g = (1/\mu_Y,\, -\mu_X/\mu_Y^2)^\top$.
Çok boyutlu delta:
\[
  \sqrt{n}(R_n - \mu_X/\mu_Y) \xrightarrow{d} \Normal\!\left(0,\,
  \frac{\sigma_X^2}{\mu_Y^2} - \frac{2\mu_X\sigma_{XY}}{\mu_Y^3} + \frac{\mu_X^2\sigma_Y^2}{\mu_Y^4}
  \right).
\]
\end{ornek}

% =====================================================================
\section{Karakteristik Fonksiyon: Tersçevirme ve Benzersizlik}
\label{sec:kf_ters}
% =====================================================================

\begin{motivasyon}[KF Neden Bu Kadar Güçlü?]
Moment üreten fonksiyon her dağılım için mevcut olmayabilir (Cauchy).
Karakteristik fonksiyon $\varphi_X(t) = E[e^{itX}]$ her dağılım için vardır.
Dahası, KF dağılımı \emph{tam olarak} belirler: iki rasgele değişkenin KF'leri özdeşse
dağılımları da özdeştir. Tersçevirme formülü, KF'den PDF'ye geçiş sağlar.
\end{motivasyon}

\begin{teorem}[KF Benzersizlik Teoremi]
\label{thm:kf_benzersizlik}
$X$ ve $Y$ rasgele değişkenler. Eğer $\varphi_X(t) = \varphi_Y(t)$ tüm $t \in \mathbb{R}$ için
geçerliyse $X \overset{d}{=} Y$ (aynı dağılım).
\end{teorem}

\begin{proof}[İspat Taslağı]
$F_X = F_Y$ olduğu yeter. Herhangi $a < b$ için
\[
  P(a < X \leq b) = \lim_{T\to\infty} \frac{1}{2\pi} \int_{-T}^T \frac{e^{-ita}-e^{-itb}}{it}\varphi_X(t)\,dt,
\]
(Lévy tersçevirme formülü). Aynı formül $F_Y$ için de geçerli; $\varphi_X = \varphi_Y$ koşuluyla
$F_X(b)-F_X(a) = F_Y(b)-F_Y(a)$ her $a < b$ için. $F_X = F_Y$ çıkar.
\end{proof}

\begin{teorem}[Tersçevirme Formülü]
\label{thm:kf_terscevir}
$X$ sürekli rasgele değişken, $\int_{-\infty}^{+\infty} |\varphi_X(t)|\,dt < \infty$ ise
\[
  f_X(x) = \frac{1}{2\pi} \int_{-\infty}^{+\infty} e^{-itx}\,\varphi_X(t)\,dt.
\]
Yani KF, PDF'nin Fourier dönüşümüdür; tersçevirme Fourier ters dönüşümüdür.
\end{teorem}

\begin{ornek}[Normal KF'sinden PDF'ye]
$\varphi_X(t) = e^{i\mu t - \sigma^2 t^2/2}$ (Normal KF).
$\int |\varphi_X(t)|\,dt = \int e^{-\sigma^2 t^2/2}\,dt = \sqrt{2\pi}/\sigma < \infty$.
Tersçevirme:
\[
  f_X(x) = \frac{1}{2\pi}\int_{-\infty}^{+\infty} e^{-itx} e^{i\mu t - \sigma^2 t^2/2}\,dt
  = \frac{1}{\sigma\sqrt{2\pi}} e^{-(x-\mu)^2/(2\sigma^2)}.
\]
(Gauss integralini tamamlayarak $t$'deki kareyi tamamlama kullanılır.)
\end{ornek}

\begin{teorem}[Lévy Süreklilik Teoremi]
\label{thm:levy_sureklilik}
$\varphi_{X_n}(t) \to \varphi(t)$ her $t$'de, ve $\varphi$ $t=0$'da sürekli ise
$X_n \xrightarrow{d} X$ burada $X$'in KF'si $\varphi$'dir.
\end{teorem}

\begin{proof}[İspat Notu]
Temel fikir: KDF dizisinin görece kompaktlığı (Helly seçim teoremi) + KF yakınsamasının
limit ölçümünü tam belirlemesi. Tam ispat için Billingsley~\cite{Billingsley1995}, §26.
Bu teorem MLT'nin KF ile ispatının kalbini oluşturur (Bölüm~\ref{chp:merkezi_limit}).
\end{proof}

\begin{hataanalizi}
"$\varphi_{X_n}(t) \to \varphi(t)$ her $t$'de, dolayısıyla $X_n \xrightarrow{d} X$."
Bu çıkarım Lévy Sürekliliği'ni atlıyor: $\varphi$'nin $t=0$'da sürekli (yani geçerli bir KF) olması
gerekir. $\varphi_n(t) = e^{-n|t|}$ ise limit $\varphi(0)=1$, $\varphi(t)=0$ ($t\neq0$) —
bu bir KF değildir. Sonuç: $X_n \to \infty$ dağılıma (mass escape), limit dağılım mevcut değil.
\end{hataanalizi}

\begin{mikroozet}
KF her dağılım için vardır; benzersiz olarak belirler.
Tersçevirme: $f_X = (2\pi)^{-1}\int e^{-itx}\varphi_X(t)\,dt$ ($|\varphi|$ integrallenebilir ise).
Lévy Sürekliliği: KF yakınsaması $\Leftrightarrow$ dağılımla yakınsama.
\end{mikroozet}

% =====================================================================
\section{İkinci Dereceden Delta Yöntemi}
\label{sec:delta2}
% =====================================================================

\begin{motivasyon}[Birinci Türev Sıfır Olduğunda]
Delta yöntemi $g'(\mu) \neq 0$ gerektirir. $g'(\mu) = 0$ ise $\sqrt{n}(g(\bar{X}_n) - g(\mu))$
sıfıra yakınsar: limit dağılım dejenere. Doğru ölçek $n$'dir; limit dağılım normal değil,
ki-kare ile ilgilidir.
\end{motivasyon}

\begin{teorem}[İkinci Dereceden Delta Yöntemi]
\label{thm:delta2}
$\sqrt{n}(X_n - \mu) \xrightarrow{d} \Normal(0, \sigma^2)$, $g'(\mu) = 0$,
$g''(\mu) \neq 0$ ise:
\[
  n(g(X_n) - g(\mu)) \xrightarrow{d} \frac{\sigma^2 g''(\mu)}{2} \chi^2_1.
\]
\end{teorem}

\begin{proof}
İkinci derecede Taylor:
$g(X_n) - g(\mu) = g'(\mu)(X_n-\mu) + \tfrac{g''(\mu)}{2}(X_n-\mu)^2 + o((X_n-\mu)^2)$.
$g'(\mu) = 0$:
$n(g(X_n)-g(\mu)) = \tfrac{g''(\mu)}{2}\cdot n(X_n-\mu)^2 + o_P(1)$.
$\sqrt{n}(X_n-\mu) \xrightarrow{d} \Normal(0,\sigma^2)$;
sürekli haritalama: $n(X_n-\mu)^2 \xrightarrow{d} \sigma^2 Z^2$ ($Z \sim \Normal(0,1)$).
$Z^2 \sim \chi^2_1$: $n(g(X_n)-g(\mu)) \xrightarrow{d} \frac{\sigma^2 g''(\mu)}{2}\chi^2_1$.
\end{proof}

\begin{ornek}[Varyans Tahmincisi]
$X_i$ i.i.d., $E[X_i] = \mu$, $\mathrm{Var}(X_i) = \sigma^2$, $E[(X_i-\mu)^4] = \mu_4 < \infty$.
$g(m) = (m - \mu)^2$ fonksiyonu, $g'(\mu) = 0$, $g''(\mu) = 2$.

$n(\bar{X}_n - \mu)^2 \xrightarrow{d} \sigma^2 \chi^2_1$.

Örneğin $\sigma^2 = 1$: $n(\bar{X}_n-\mu)^2 \xrightarrow{d} \chi^2_1$.
Bu, varyans-örneklem büyüklüğü ilişkisinin ikinci dereceden hassas analizini sağlar.
\end{ornek}

\begin{ornek}[Sıfır Noktasında Sinüs Dönüşümü]
$\bar{X}_n \xrightarrow{P} 0$; $g(x) = \sin(x)$. $g'(0) = 1 \neq 0$, normal delta yöntemi geçerli.
Ama $h(x) = 1 - \cos(x)$: $h'(0) = 0$, $h''(0) = 1$.
$n(1 - \cos(\bar{X}_n)) \xrightarrow{d} \frac{\sigma^2}{2}\chi^2_1$.
\end{ornek}

% =====================================================================
\section{Olasılık İntegral Dönüşümü ve Kuantil Dönüşümü}
\label{sec:prob_integral}
% =====================================================================

\begin{teorem}[Olasılık İntegral Dönüşümü (OİD)]\label{thm:oid}
$X$ sürekli rasgele değişken, KDF $F$ kesinlikle artan ve sürekli. O zaman
\[
  U = F(X) \sim \Uniform(0,1).
\]
\end{teorem}

\begin{proof}
$P(U \leq u) = P(F(X)\leq u) = P(X\leq F^{-1}(u)) = F(F^{-1}(u)) = u$, $u\in(0,1)$. Bu, $\Uniform(0,1)$ KDF'sidir.
\end{proof}

\begin{teorem}[Kuantil dönüşümü]\label{thm:kuantil}
$U\sim\Uniform(0,1)$ ve $F$ herhangi bir KDF, $F^{-1}$ genelleştirilmiş tersi. O zaman $X=F^{-1}(U)$ dağılımı $F$'dir.
\end{teorem}

\begin{proof}
$P(X\leq x)=P(F^{-1}(U)\leq x)=P(U\leq F(x))=F(x)$.
\end{proof}

\begin{ornek}[Simülasyon yöntemi]\label{ex:simulasyon}
$\Exp(\lambda)$ simülasyonu: $F(x)=1-e^{-\lambda x}$, $F^{-1}(u)=-\log(1-u)/\lambda$. $U\sim\Uniform(0,1)$ üret; $X=-\log U/\lambda$ hesapla ($1-U\overset{d}{=}U$).
Box-Muller ($\Normal(0,1)$): $U_1,U_2\iid\Uniform(0,1)$:
\[
  Z_1=\sqrt{-2\log U_1}\cos(2\pi U_2),\quad Z_2=\sqrt{-2\log U_1}\sin(2\pi U_2)\iid\Normal(0,1).
\]
\end{ornek}

\begin{hataanalizi}
OİD yalnızca sürekli dağılımlar için geçerlidir. Kesikli $X$ için $F(X)$ $\Uniform(0,1)$'e dağılmaz; bunun yerine $U\in(F(X^-),F(X)]$ aralığında üretilmiş $\Uniform$ ile çalışmak gerekir.
\end{hataanalizi}

% =====================================================================
\section{Çarpışma Transformasyonları ve Simetri}
\label{sec:simetri_donusumler}
% =====================================================================

\begin{tanim}[Simetrik dağılım]\label{def:simetrik}
$X$ dağılımı $0$ çevresinde simetriktir: $X\overset{d}{=}-X$.
\end{tanim}

\begin{teorem}[Mutlak değer dönüşümü]\label{thm:mutlak}
$X$, $0$ çevresinde simetrik, $f$ yoğunluğu mevcut. $Y=|X|$'in yoğunluğu:
\[
  f_Y(y) = 2f(y),\quad y>0.
\]
\end{teorem}

\begin{proof}
$P(|X|\leq y)=P(-y\leq X\leq y)=F(y)-F(-y)$. Türev: $f_Y(y)=f(y)+f(-y)=2f(y)$ (simetri).
\end{proof}

\begin{ornek}[Yarı-normal]\label{ex:yari_normal}
$Z\sim\Normal(0,1)$: $|Z|$ yarı-normal (\emph{half-normal}), $f_{|Z|}(y)=2\phi(y)$, $y>0$.
$\Beklenti[|Z|]=\sqrt{2/\pi}$, $\Var(|Z|)=1-2/\pi$.
\end{ornek}

\begin{teorem}[İşaret-büyüklük bağımsızlığı]\label{thm:isaret_buyukluk}
$X$ $0$ çevresinde simetrik. $S=\mathrm{sgn}(X)$ ve $M=|X|$ bağımsızdır; $S\sim\mathrm{Bernoulli}(1/2)$.
\end{teorem}

\begin{proof}
$P(S=1,M\leq m)=P(0<X\leq m)=F(m)-1/2=\tfrac12 P(M\leq m)$. Benzer şekilde $P(S=-1,M\leq m)=\tfrac12 P(M\leq m)$: bağımsızlık.
\end{proof}

\begin{disiplinlerarasibaglan}
\textbf{Finansal Matematik:} OİD, Monte Carlo risk simülasyonunun temelidir. Sklar teoremi: her ortak dağılım $F(\mathbf{x})=C(F_1(x_1),\ldots,F_p(x_p))$ biçiminde yazılır; $C$ \emph{copula}'dır — CDO kredi riski modellemesinde kullanılır.

\textbf{Kalite Mühendisliği:} Yarı-normal dağılım, Taguchi yönteminde tek taraflı üretim toleranslarını modellemek için kullanılır.
\end{disiplinlerarasibaglan}

\begin{mikroozet}
$g'(\mu) = 0$: standart delta yöntemi çöker.
İkinci derecede Delta: $n(g(X_n)-g(\mu)) \to \frac{\sigma^2 g''(\mu)}{2}\chi^2_1$.
Limit dağılım ki-kare; ölçek $\sqrt{n}$ yerine $n$.
OİD: $F(X)\sim\Uniform(0,1)$ — simülasyon ve copula modellemesinin temeli.
Simetrik dağılım: işaret ve büyüklük bağımsız; yarı-normal $\Beklenti[|Z|]=\sqrt{2/\pi}$.
\end{mikroozet}

% ---- Disiplinlerarası Bağlantı (TYMM) ----
\begin{kopru}[Disiplinlerarası — Mühendislik ve Veri Bilimi]
\textbf{Güvenilirlik mühendisliği:} Paralel sistemlerde hayatta kalma süresi $X_{(n)}$,
seri sistemlerde $X_{(1)}$. "$n$ bileşenden $k$'si çalışırsa sistem çalışır"
modeli, $k$'ıncı sıra istatistiğinin dağılımına dayanır.\\
\textbf{Veri bilimi:} Örneklem medyanı, aşırı değerlere dirençli bir merkez ölçüsüdür;
sıra istatistiklerinden türetilen yüzdelikler ($Q_1$, $Q_3$) kutu grafiğinin temelidir.\\
\textbf{Ekonometri:} Delta yöntemi, doğrusal olmayan parametre dönüşümlerinde
güven aralığı hesabında standart araçtır (oran tahmincisi, elastisite).
\end{kopru}

% ---- Öz-Yansıtma ----
\begin{ozyansima}
\begin{itemize}[leftmargin=1.8em]
  \item MÜF yönteminin Jacobian yöntemine göre avantajı nedir? Hangi durumda
        Jacobian daha kullanışlıdır?
  \item Sıra istatistiklerinin neden bağımlı olduğunu kendi cümlelerinizle açıklayın.
  \item Delta yönteminin temel fikri nedir? Taylor açılımının bu bağlamda nasıl çalıştığını anlatın.
  \item Üstel dağılımlı cihazlarda sıra arası boşlukların bağımsız olması sezgiye uygun mu?
\end{itemize}
\end{ozyansima}

% ---- Köprü Notu ----
\begin{kopru}
\textbf{Önceki bölüm:} Bölüm~\ref{chp:dagilim_aileleri} — $t$, $F$, $\chi^2$ dağılımları
burada türetildi; Bölüm~\ref{chp:beklenti}'deki MÜF/KF araçları evrişimde kullanıldı.\\
\textbf{Sonraki bölüm:} Bölüm~\ref{chp:merkezi_limit} — Delta yöntemi ve KF
yöntemi Merkezi Limit Teoremi'nin ispatında merkezi rol oynayacak.\\
\textbf{İstatistik kısmı:} Sıra istatistikleri Böl.~\ref{chp:istatistik_temelleri}'de
yüzdelik tahmincisi ve Böl.~\ref{chp:hipotez_testi}'nde parametresiz testlerde kullanılır.
\defterbaglantisi{bolum06\_donusumler.ipynb}{%
  Jacobian görselleştirmesi, sıra istatistikleri simülasyonu, delta yöntemi uygulaması}
\end{kopru}

% =====================================================================
%  ALIŞTIRMALAR
% =====================================================================

\cozumlualistirmalar

\begin{alistirma}[Al.6.1]{A}{Tekil Dönüşüm}
$X \sim \mathrm{Exp}(1)$. $Y = \sqrt{X}$ için:
(a) $f_Y(y)$'yi Jacobian yöntemiyle bulun.
(b) $E[Y]$'yi $f_Y$ üzerinden hesaplayın ($\Gamma(3/2) = \sqrt{\pi}/2$).
\end{alistirma}
\alcozum{%
(a) $g(x) = \sqrt{x}$, ters $h(y) = y^2$, $|h'(y)| = 2y$.
$f_Y(y) = e^{-y^2} \cdot 2y$, $y > 0$ (Rayleigh dağılımı, Bölüm~\ref{chp:rasgele_degiskenler}'de de elde edilmişti).\\
(b) $E[Y] = \int_0^\infty y \cdot 2y e^{-y^2} dy = 2\int_0^\infty y^2 e^{-y^2} dy$.
$u = y^2$: $= \int_0^\infty u^{1/2} e^{-u} du = \Gamma(3/2) = \sqrt{\pi}/2$.}

\begin{alistirma}[Al.6.2]{A}{Sıra İstatistikleri}
$X_1, X_2, X_3 \sim \Uniform(0,1)$ bağımsız.
(a) $X_{(1)}$ ve $X_{(3)}$'ün PDF'lerini yazın.
(b) $E[X_{(1)}]$ ve $E[X_{(3)}]$'ü hesaplayın.
(c) $P(X_{(2)} > 0.5)$'i bulun.
\end{alistirma}
\alcozum{%
(a) $f_{X_{(1)}}(x) = 3(1-x)^2$, $f_{X_{(3)}}(x) = 3x^2$, $0<x<1$.\\
(b) $E[X_{(1)}] = 1/(3+1) = 1/4$. $E[X_{(3)}] = 3/4$.\\
(c) $X_{(2)} \sim \Beta(2,2)$: $f_{X_{(2)}}(x) = 6x(1-x)$.
$P(X_{(2)}>0.5) = \int_{0.5}^1 6x(1-x)dx = 6[x^2/2-x^3/3]_{0.5}^1
= 6[(1/2-1/3)-(1/8-1/24)] = 6[1/6 - 1/12] = 6 \cdot 1/12 = 1/2$. (Simetri!)}

\begin{alistirma}[Al.6.3]{B}{MÜF Yöntemi}
$X_1, \ldots, X_n \sim \Pois(\lambda)$ bağımsız. $S_n = \sum X_i$.
(a) $M_{S_n}(t)$'yi bulun.
(b) $S_n$'nin dağılımını belirtin.
(c) $\bar{X}_n = S_n/n$'nin MÜF'ünü yazın ve $n=100$, $\lambda=2$ için
$P(\bar{X}_{100} \geq 2.5)$'e Chernoff sınırı verin.
\end{alistirma}
\alcozum{%
(a) $M_{S_n}(t) = \prod M_{X_i}(t) = [e^{\lambda(e^t-1)}]^n = e^{n\lambda(e^t-1)}$.\\
(b) $S_n \sim \Pois(n\lambda)$.\\
(c) $M_{\bar{X}_n}(t) = E[e^{t\sum X_i/n}] = M_{S_n}(t/n) = e^{n\lambda(e^{t/n}-1)}$.
Chernoff: $\Lambda^*(a) = \sup_t(ta - n\lambda(e^{t/n}-1))/n$. Optimize: $t^* = n\ln(a/\lambda)$.
$\Lambda^*(2.5) = 2.5\ln(2.5/2) - 2(2.5/2 - 1) = 2.5\ln(1.25) - 0.5 \approx 0.558 - 0.5 = 0.058$.
$P(\bar{X}_{100}\geq2.5) \leq e^{-100\cdot0.058} = e^{-5.8} \approx 0.003$.}

\begin{alistirma}[Al.6.4]{B}{Delta Yöntemi}
$X_1, \ldots, X_n$ i.i.d., $E[X_i] = \mu > 0$, $\mathrm{Var}(X_i) = \sigma^2$.
MLT (bir sonraki bölüm): $\sqrt{n}(\bar{X}_n - \mu) \xrightarrow{d} \Normal(0,\sigma^2)$.
Delta yöntemi ile $\sqrt{n}(1/\bar{X}_n - 1/\mu)$'nun limit dağılımını bulun.
\end{alistirma}
\alcozum{%
$g(x) = 1/x$, $g'(x) = -1/x^2$, $g'(\mu) = -1/\mu^2$.
Delta yöntemi: $\sqrt{n}(g(\bar{X}_n)-g(\mu)) \xrightarrow{d} \Normal(0, \sigma^2/\mu^4)$.
$\sqrt{n}(1/\bar{X}_n - 1/\mu) \xrightarrow{d} \Normal(0, \sigma^2/\mu^4)$.}

\begin{alistirma}[Al.6.4b]{B}{İki Boyutlu Jacobian — Polar Dönüşüm}
$(X,Y)\iid\Normal(0,1)$ bağımsız. $R=\sqrt{X^2+Y^2}$, $\Theta=\arctan(Y/X)$ polar dönüşümü.
(a) $(R,\Theta)$'nın birleşik yoğunluğunu Jacobian yöntemiyle bulun.
(b) $R$ ve $\Theta$'nın marjinal dağılımlarını yazın. $R$'nin adı ne?
(c) Bu sonucu Box-Muller dönüşümüyle ilişkilendirin.
\end{alistirma}
\alcozum{%
(a) Ters dönüşüm: $X=R\cos\Theta$, $Y=R\sin\Theta$. Jacobian:
\[
  J = \begin{vmatrix}\cos\Theta & -R\sin\Theta \\ \sin\Theta & R\cos\Theta\end{vmatrix} = R.
\]
Birleşik yoğunluk: $f_{X,Y}(x,y)=\frac{1}{2\pi}e^{-(x^2+y^2)/2}$.
$f_{R,\Theta}(r,\theta) = f_{X,Y}(r\cos\theta, r\sin\theta)\cdot|J|^{-1} = \frac{1}{2\pi}e^{-r^2/2}\cdot r$, $r>0$, $\theta\in[0,2\pi)$.

(b) $f_\Theta(\theta)=\frac{1}{2\pi}$ (Uniform$[0,2\pi)$); $f_R(r)=re^{-r^2/2}$ ($r>0$): Rayleigh dağılımı. $R^2\sim\chi^2_2=\mathrm{Exp}(1/2)$.

(c) Box-Muller: $U_1,U_2\iid\Uniform[0,1]$'den $X=\sqrt{-2\log U_1}\cos(2\pi U_2)$, $Y=\sqrt{-2\log U_1}\sin(2\pi U_2)$ bağımsız standart normaldir. Burada $R=\sqrt{-2\log U_1}\sim$ Rayleigh ve $\Theta=2\pi U_2\sim\Uniform[0,2\pi)$; bu tam olarak $(R,\Theta)$'nın dağılımıdır.}

\begin{alistirma}[Al.6.4c]{C}{Konvolüsyon ve Toplam Dağılımı}
$X\sim\Gama(\alpha,1)$, $Y\sim\Gama(\beta,1)$ bağımsız.
(a) MÜF yöntemiyle $Z=X+Y$'nin dağılımını bulun.
(b) $W=X/(X+Y)$'nin dağılımını konvolüsyon + Jacobian ile türetin.
(c) $\alpha=\beta=1$ durumunda (yani $X,Y\iid\mathrm{Exp}(1)$): $W$'nin dağılımını yazın ve sezgisel yorumunu verin.
\end{alistirma}
\alcozum{%
(a) $M_X(t)=(1-t)^{-\alpha}$, $M_Y(t)=(1-t)^{-\beta}$ (t<1). $M_Z(t)=(1-t)^{-(\alpha+\beta)}$; dolayısıyla $Z\sim\Gama(\alpha+\beta,1)$.

(b) Değişken dönüşümü: $(X,Y)\mapsto(Z=X+Y,W=X/(X+Y))$. Ters: $X=WZ$, $Y=(1-W)Z$. Jacobian: $|J|=Z$. Birleşik: $f_{Z,W}(z,w)=f_X(wz)f_Y((1-w)z)\cdot z$.
\[
  f_X(x)f_Y(y) = \frac{x^{\alpha-1}e^{-x}}{\Gamma(\alpha)}\cdot\frac{y^{\beta-1}e^{-y}}{\Gamma(\beta)}.
\]
$f_{Z,W}(z,w)=\frac{(wz)^{\alpha-1}((1-w)z)^{\beta-1}e^{-z}}{\Gamma(\alpha)\Gamma(\beta)}\cdot z = \frac{z^{\alpha+\beta-1}e^{-z}}{\Gamma(\alpha+\beta)}\cdot\frac{w^{\alpha-1}(1-w)^{\beta-1}}{B(\alpha,\beta)}$.

Marjinal $w$: $f_W(w)=\frac{w^{\alpha-1}(1-w)^{\beta-1}}{B(\alpha,\beta)}$; yani $W\sim\Beta(\alpha,\beta)$.

(c) $\alpha=\beta=1$: $W\sim\Beta(1,1)=\Uniform[0,1]$. Yorum: İki bağımsız Üstel sürenin toplamındaki birincinin oranı, tüm $[0,1]$'de eşit olasılıklı. Yenileme teorisi: bekleme kesiri üniform.}

% ---

\alistirmalar

\begin{alistirma}[Al.6.5]{A}{Dönüşüm Uygulaması}
$X \sim \Normal(0,1)$. $Y = e^X$ (lognormal).
(a) $F_Y(y)$'yi $\Phi$ cinsinden yazın. (b) $f_Y(y)$'yi türetin. (c) $P(1 < Y < e)$'yi bulun.
\end{alistirma}
\alcozum{%
(a) $F_Y(y) = P(e^X \leq y) = P(X \leq \ln y) = \Phi(\ln y)$.
(b) $f_Y(y) = \phi(\ln y)/y = e^{-(\ln y)^2/2}/(y\sqrt{2\pi})$, $y>0$.
(c) $P(1<Y<e) = P(0<X<1) = \Phi(1)-\Phi(0) \approx 0.8413 - 0.5 = 0.3413$.}

\begin{alistirma}[Al.6.6]{A}{Minimum Dağılımı}
$X_1,\ldots,X_5 \sim \mathrm{Exp}(2)$ bağımsız. $X_{(1)} = \min X_i$.
(a) $X_{(1)}$'in dağılımını bulun. (b) $P(X_{(1)} > 0.3)$'ü hesaplayın.
(c) $E[X_{(1)}]$'i bulun.
\end{alistirma}
\alcozum{%
(a) $X_{(1)} \sim \mathrm{Exp}(5\cdot2) = \mathrm{Exp}(10)$.
(b) $P(X_{(1)}>0.3) = e^{-10\cdot0.3} = e^{-3} \approx 0.050$.
(c) $E[X_{(1)}] = 1/10 = 0.1$.}

\begin{alistirma}[Al.6.7]{B}{Çok Boyutlu Jacobian}
$X, Y$ bağımsız $\Normal(0,1)$. $U = X+Y$, $V = X-Y$.
(a) $(U,V)$'nin ortak PDF'sini bulun.
(b) $U$ ve $V$ bağımsız mı? Dağılımlarını belirtin.
(c) $U$ ve $V$'nin $\mathrm{Cov}(U,V)$'sini hesaplayın.
\end{alistirma}
\alcozum{%
(a) Ters: $x=(u+v)/2$, $y=(u-v)/2$. $|\det J| = 1/2$.
$f_{U,V}(u,v) = \phi((u+v)/2)\phi((u-v)/2)/2 = \frac{1}{4\pi}e^{-(u^2+v^2)/4}$.
(b) $f_{U,V} = f_U(u)f_V(v)$ ile $U \sim \Normal(0,2)$, $V \sim \Normal(0,2)$, bağımsız.
(c) $\mathrm{Cov}(U,V) = \mathrm{Cov}(X+Y,X-Y) = \mathrm{Var}(X)-\mathrm{Var}(Y) = 1-1 = 0$. (Bağımsızlıkla tutarlı.)}

\begin{alistirma}[Al.6.8]{B}{Sıra İstatistikleri — Beklenti}
$X_1,\ldots,X_n \sim \Uniform(0,\theta)$. $X_{(n)} = \max X_i$ ile $\theta$'yı tahmin edin.
(a) $X_{(n)}$'nin PDF'sini bulun. (b) $E[X_{(n)}]$'i hesaplayın.
(c) $E[X_{(n)}] \neq \theta$: yansız tahminci nedir?
\end{alistirma}
\alcozum{%
(a) $F(x) = x/\theta$, $f_{X_{(n)}}(x) = n(x/\theta)^{n-1}/\theta = nx^{n-1}/\theta^n$.
(b) $E[X_{(n)}] = \int_0^\theta x \cdot nx^{n-1}/\theta^n dx = n\theta/(n+1)$.
(c) $\hat\theta = (n+1)/n \cdot X_{(n)}$ yansız tahminci: $E[\hat\theta] = \theta$.}

\begin{alistirma}[Al.6.9]{C}{Evrişim: Gama Toplamı}
$X \sim \Gama(\alpha, \beta)$, $Y \sim \Gama(\gamma, \beta)$ bağımsız ($\beta$ aynı).
Evrişim integraliyle (MÜF kullanmadan) $X+Y \sim \Gama(\alpha+\gamma, \beta)$
olduğunu ispatlayın.
\end{alistirma}
\alcozum{%
$f_{X+Y}(z) = \int_0^z f_X(t)f_Y(z-t)dt
= \int_0^z \frac{\beta^\alpha t^{\alpha-1}e^{-\beta t}}{\Gamma(\alpha)} \cdot \frac{\beta^\gamma(z-t)^{\gamma-1}e^{-\beta(z-t)}}{\Gamma(\gamma)} dt$.
$= \frac{\beta^{\alpha+\gamma}e^{-\beta z}}{\Gamma(\alpha)\Gamma(\gamma)} \int_0^z t^{\alpha-1}(z-t)^{\gamma-1}dt$.
$u = t/z$: $\int_0^z t^{\alpha-1}(z-t)^{\gamma-1}dt = z^{\alpha+\gamma-1}B(\alpha,\gamma)
= z^{\alpha+\gamma-1}\Gamma(\alpha)\Gamma(\gamma)/\Gamma(\alpha+\gamma)$.
$f_{X+Y}(z) = \frac{\beta^{\alpha+\gamma}}{\Gamma(\alpha+\gamma)}z^{\alpha+\gamma-1}e^{-\beta z}$
— bu $\Gama(\alpha+\gamma,\beta)$ PDF'sidir.}

\begin{alistirma}[Al.6.10]{B}{KF Tersçevirme}
$X$ rasgele değişken, $\varphi_X(t) = e^{-|t|}$ (Cauchy KF).
(a) $\int_{-\infty}^{+\infty} |\varphi_X(t)|\,dt$'nin varlığını kontrol edin.
(b) Tersçevirme formülünü kullanarak $f_X(x)$'i hesaplayın.
(c) Sonucun hangi standart dağılıma karşılık geldiğini belirtin.
\end{alistirma}
\alcozum{%
(a) $\int_{-\infty}^\infty |e^{-|t|}|\,dt = 2\int_0^\infty e^{-t}\,dt = 2 < \infty$. Formül uygulanabilir.\\
(b) $f_X(x) = \frac{1}{2\pi}\int_{-\infty}^{\infty} e^{-itx}e^{-|t|}\,dt
= \frac{1}{2\pi}\left[\int_0^\infty e^{t(-1-ix)}\,dt + \int_0^\infty e^{t(-1+ix)}\,dt\right]$
$= \frac{1}{2\pi}\left[\frac{1}{1+ix} + \frac{1}{1-ix}\right] = \frac{1}{2\pi}\cdot\frac{2}{1+x^2} = \frac{1}{\pi(1+x^2)}$.\\
(c) $f_X(x) = \frac{1}{\pi(1+x^2)}$ — standart Cauchy$(0,1)$ PDF'si.}

\begin{alistirma}[Al.6.11]{B}{İkinci Dereceden Delta Yöntemi}
$X_1, \ldots, X_n$ i.i.d., $E[X_i] = 0$, $\mathrm{Var}(X_i) = 1$.
$g(x) = x^2 + 3x$.
(a) $g'(0)$ ve $g''(0)$'ı hesaplayın.
(b) $\sqrt{n}(g(\bar{X}_n) - g(0))$ için limit dağılımı verin.
(c) Eğer $h(x) = x^2$ olsaydı, $n(h(\bar{X}_n) - h(0))$ için limit dağılımını verin.
\end{alistirma}
\alcozum{%
$g(x) = x^2 + 3x$; $g'(x) = 2x+3$; $g'(0) = 3 \neq 0$; $g''(0) = 2$.\\
(a) $g'(0) = 3$, $g''(0) = 2$.\\
(b) Standart delta (çünkü $g'(0)=3\neq0$): $\sqrt{n}(g(\bar{X}_n)-0) \xrightarrow{d} \Normal(0, 1\cdot 9) = \Normal(0,9)$.\\
(c) $h(x)=x^2$; $h'(0)=0$, $h''(0)=2$. İkinci dereceden delta:
$n(h(\bar{X}_n)-0) \xrightarrow{d} \frac{1\cdot2}{2}\chi^2_1 = \chi^2_1$.
Yani $n\bar{X}_n^2 \xrightarrow{d} \chi^2_1$ — bu MLT'nin doğal bir sonucudur.}

% =====================================================================
\endinput
